## A stronger—and materially more robust—candidate

The revised argument now gives:

> **Strengthened candidate theorem.**  
> Let \(X\) be a primitive distance-regular graph on \(n\) vertices with diameter \(d\ge 3\). Then either \(X\) is a Johnson graph or a Hamming graph, or
> \[
> \boxed{\operatorname{motion}(X)\ge \frac{n}{8d^{3}}.}
> \]

This improves the previous candidate

\[
\frac{n}{32d^{3}}
\]

by a factor of \(4\).

There is also a deliberately more conservative version:

> **Source-conservative candidate.**  
> Using only the published structural inputs from Pyber–Skresanov and Kivva,
> \[
> \boxed{\operatorname{motion}(X)\ge \frac{n}{15d^{3}}.}
> \]

The \(1/(8d^3)\) version additionally uses a January 2026 preprint of Lv and Koolen, but only in the narrow case \(d=3\) and \(\mu=c_2\ge3\). Their structural theorem classifies the relevant geometric possibilities into Johnson, Grassmann, locally disconnected, and two residual parameter configurations; the inequalities obtained from small motion eliminate every case except Johnson. citeturn417751view3turn417751view4

Both statements remain **machine-generated candidate theorems**, not established results. The published benchmark is \(C n/d^6\), outside the Johnson, Hamming, and crown families, so even the conservative version would improve the known diameter exponent from \(6\) to \(3\). citeturn341760view0turn341760view1turn341760view3

## The most robust breakthrough: an exact geodesic Poincaré inequality

The strongest standalone contribution is independent of the delicate Johnson/Hamming classification.

> **Geodesic spectral-gap lemma.**  
> Let \(G\) be the graph of a connected symmetric basis relation of valency \(k\) in a homogeneous coherent configuration on \(n\) points. If \(d\) is its diameter and \(\theta_1\) its second adjacency eigenvalue, then
> \[
> \boxed{
> k-\theta_1\ge
> \frac{n^2k}
> {\displaystyle\sum_{x,y\in V(G)}
> \operatorname{dist}(x,y)^2}
> \ge\frac{k}{d^2}.
> }
> \]

For a distance-regular graph this specializes to

\[
\boxed{
k-\theta_1\ge
\frac{nk}{\displaystyle\sum_{i=0}^{d}k_i i^2}.
}
\]

Pyber and Skresanov proved

\[
k-\theta_1\ge\frac{k}{8d^2}
\]

by first obtaining edge expansion \(k/(2d)\) from uniform geodesic loads and then applying Cheeger’s inequality. The new proof uses the same uniformity phenomenon but applies Cauchy–Schwarz directly along the geodesics, retaining their lengths rather than discarding them. This removes the factor \(8\). citeturn417751view1turn341760view2

### Proof in one page

For every ordered pair \(x,y\), let \(p(x,y)\) be the number of geodesics from \(x\) to \(y\). For a directed edge \(e\), define

\[
Q_e=
\sum_{x,y}
\frac{\operatorname{dist}(x,y)}{p(x,y)}
\#\{P:P\text{ is an }x\text{--}y
\text{ geodesic containing }e\}.
\]

The coherent-configuration intersection-number identities imply that \(Q_e\) is independent of the directed edge \(e\); call the common value \(Q\). Summing over all \(nk\) directed edges gives

\[
nkQ=\sum_{x,y}\operatorname{dist}(x,y)^2.
\]

For a mean-zero function \(f\), Cauchy–Schwarz along each geodesic gives

\[
(f(x)-f(y))^2
\le
\frac{\operatorname{dist}(x,y)}{p(x,y)}
\sum_{P:x\to y}\sum_{e\in P}(\nabla_e f)^2.
\]

Summing over all ordered \(x,y\),

\[
2n\lVert f\rVert_2^2
\le
Q\sum_{e\text{ directed}}(\nabla_e f)^2
=
2Q\,f^{\mathsf T}(kI-A)f.
\]

Taking \(f\) in the \(\theta_1\)-eigenspace yields

\[
k-\theta_1\ge\frac nQ
=
\frac{n^2k}{\sum_{x,y}\operatorname{dist}(x,y)^2}.
\]

I did not locate this exact coherent-configuration formulation in the literature I checked. It may be subsumed by a known canonical-path or multicommodity-flow inequality, so I am not claiming novelty without a broader search. The proof itself is short and comparatively easy to audit.

## A cleaner small-support structural lemma

The argument now begins with an exact identity for the number \(D(1)\) of vertices distinguishing an adjacent pair:

\[
\boxed{
D(1)=2+\frac2k\sum_{i=2}^{d}k_i c_i.
}
\]

Consequently,

\[
\boxed{
D(1)>\frac{\mu}{k}n,
\qquad \mu=c_2.
}
\]

Now suppose an automorphism \(g\) has support \(S\), with density

\[
\rho=\frac{|S|}{n}\le\frac12.
\]

The support-sensitive final inequality in Pyber–Skresanov’s geodesic expansion proof gives

\[
|\delta(S)|
\ge
|S|\frac{k}{d}(1-\rho).
\]

Thus some moved vertex \(x\) has at least \(k(1-\rho)/d\) fixed neighbors. In the relevant small-\(\mu\) regime, those fixed neighbors force \(x\) and \(x^g\) to be adjacent. Since every fixed vertex is equidistant from \(x\) and \(x^g\),

\[
D(1)\le |S|.
\]

It follows immediately that

\[
\mu<\rho k,
\qquad
\lambda>\frac{1-\rho}{d}k.
\]

This is a particularly useful conversion:

\[
\text{small automorphism support}
\quad\Longrightarrow\quad
\text{small }\mu\text{ and large }\lambda.
\]

Keeping the **full** clique-size expression from Metsch’s theorem, rather than replacing it by the coarser \(L\ge\lambda/2\), then gives a clique with

\[
L-1>\frac{k}{d+1}.
\]

The Delsarte clique bound forces the absolute value \(m\) of the smallest eigenvalue to satisfy

\[
m<d+1.
\]

Once the Bang–Koolen criterion produces Delsarte geometry, \(m\) is integral, so

\[
\boxed{m\le d.}
\]

The published ingredients are exactly the full Metsch expression, the Delsarte clique bound, and the Bang–Koolen condition \(m^2\mu<\lambda\). citeturn417751view0

## A sharper Hamming-stability argument

The remaining difficult branch is

\[
\mu=2,\qquad
\theta_1\ge(1-\varepsilon)b_1.
\]

Kivva’s published proof uses a sufficiently small \(\varepsilon\) on the scale

\[
\frac{1}{m^4d}.
\]

That threshold is one reason the published motion argument loses several powers of \(d\). citeturn417751view2

The new argument retains the exact standard-sequence recurrence

\[
u_{i+1}
=
u_i-\frac{k-\theta}{b_i}u_i
-\frac{c_i}{b_i}(u_{i-1}-u_i)
\]

rather than repeatedly applying a uniform worst-case loss. With

\[
A=\frac{1+(3m-1)\varepsilon}{1-m\varepsilon},
\qquad
\delta=A-1,
\]

one obtains

\[
u_{i+1}\ge
u_i\left(1-\frac{A}{m-\tau_i}\right).
\]

The key factorization is

\[
1-\frac{A}{r}
=
\frac{r-1}{r}
\left(1-\frac{\delta}{r-1}\right).
\]

The first factor telescopes; the second accumulates only harmonically. The result is

\[
\boxed{
u_j\ge
\left(1-\delta H_{j-1}\right)
\frac{m-\tau_j}
     {m-\tau_j+j-1}
\frac{\theta}{k}.
}
\]

The previous linear-in-\(d\) accumulation is thereby replaced by \(H_d=O(\log d)\).

A second exact estimate improves the size of the dominant distance sphere from the earlier coarse bound to

\[
\boxed{k_t\ge(1-2m\varepsilon)n.}
\]

Together these give enough multiplicity to force the Hamming endpoint for

\[
\varepsilon=\frac{1}{4d^3-1}
\]

in the \(1/(8d^3)\) theorem, and for

\[
\varepsilon=\frac{2}{15d^3-2}
\]

in the source-conservative theorem.

This is still the **least robust portion** of the candidate proof—not because any scalar inequality failed, but because it compresses and modifies several interacting pieces of Kivva’s multiplicity argument.

## What was verified

The accompanying audit uses exact rational arithmetic, not floating point. It checks:

- all structural inequalities for \(3\le d\le10{,}000\);
- the \(1/(8d^3)\) and \(1/(15d^3)\) parameter choices;
- every exceptional \(d=3,4\) Hamming product case;
- the exact \(d=5,6\) losses;
- the endpoint inequalities in the multiplicity factor;
- the polynomial inequalities used for the infinite \(d\ge7\) tail.

Its output is:

```text
All strengthened-candidate audit checks passed.
Structural and threshold checks: 3 <= d <= 10000
(exact rational arithmetic).
Hamming finite products and finite losses:
exact rational arithmetic.
Analytic tail inequalities:
boundary polynomials and 7 <= d <= 1000 checked.
```

The checks establish that the written numerical implications are internally consistent. They do **not** replace verification that every imported theorem has precisely the required hypotheses.

## Current confidence assessment

I would separate the result into four layers:

1. **Geodesic Poincaré inequality:** high confidence. It is self-contained and short enough for direct line-by-line checking.

2. **Adjacent distinguishing identity and small-support reduction:** high confidence.

3. **Full-Metsch argument giving \(m\le d\):** moderately high confidence. The strict inequalities and ceiling term have been retained explicitly.

4. **Hamming multiplicity argument:** moderate confidence. This is the first part that should be sent to a specialist. The exact recurrence is convincing, but the interface with Kivva’s local-eigenvalue and endpoint machinery deserves independent scrutiny.

The \(1/(15d^3)\) version is institutionally more robust because it avoids relying on a very recent preprint. Kivva explicitly records the inequality \(1/200<\varepsilon^\ast\) used by that version. citeturn417751view2

## Why the proof still stops at \(d^{-3}\)

The obstruction is now sharply identifiable. Small support density \(\rho\) gives approximately

\[
\mu\lesssim\rho k,
\qquad
m=O(d),
\qquad
\lambda\gtrsim\frac{k}{d}.
\]

The Bang–Koolen geometricity condition is

\[
m^2\mu<\lambda.
\]

Substitution gives

\[
d^2\rho k\lesssim\frac{k}{d},
\]

or

\[
\rho=O(d^{-3}).
\]

So improving the exponent to \(d^{-2}\) probably requires a new geometricity theorem that uses more of the clique structure than the scalar condition \(m^2\mu<\lambda\), or an argument that bypasses exact Delsarte geometry.

## Updated files

[Strengthened candidate proof — PDF](sandbox:/mnt/data/babai_motion_d3_strengthened_candidate.pdf)

[LaTeX source](sandbox:/mnt/data/babai_motion_d3_strengthened_candidate.tex)

[Exact-arithmetic audit script](sandbox:/mnt/data/babai_motion_d3_strengthened_audit.py)

The most defensible item to circulate first is the geodesic Poincaré proposition by itself. The full \(1/(8d^3)\) statement should be labeled a candidate until the \(\mu=2\) multiplicity section receives independent specialist review.